\subsection{Volatility, liquidity and crashes}

\begin{qfdefinition}[Fundamental and transitory volatility: the random-walk and transitory components of prices]\label{def:volatility-components}
\qffield{Statement} \emph{Volatility} is the tendency of prices to change unexpectedly. \emph{Fundamental volatility} comes from unanticipated changes in instrument values; \emph{transitory volatility} comes from trading by impatient uninformed traders, which pushes prices away from values until they revert. Total volatility is the sum of the two \citep[chap.~``Volatility'', pp.~410, 413--414]{alpha-harris}. In statistical form, the trade price is a random-walk efficient price plus a stationary pricing error \citep[sec.~8.6, p.~72]{s6-hasbrouck}:
\[ p_t=m_t+s_t, \qquad m_t=m_{t-1}+w_t, \qquad \Delta p_t=w_t+\Delta s_t, \]
where $s_t$ has mean zero, is covariance-stationary, and may be serially correlated and correlated with $w_t$. Fundamental volatility is the variance $\sigma_w^2$ of $w_t$. If the price changes have the moving-average representation $\Delta p_t=\theta(L)\varepsilon_t$, then $\sigma_w^2=\theta(1)^2\sigma_\varepsilon^2$, the same for every structural model consistent with the data, while the variance of the pricing error has only a lower bound \citep[same sec., eqs.~(8.10)--(8.13), pp.~73--74]{s6-hasbrouck}. Over long horizons $\sigma_w^2\approx\operatorname{Var}(p_t-p_{t-k})/k$, because microstructure effects matter less as the interval grows \citep[sec.~8.3, p.~69]{s6-hasbrouck}.

\qffield{Notation} $p_t$ is the trade price, $m_t$ the efficient (fundamental) price, $s_t=p_t-m_t$ the pricing error, $w_t$ the efficient-price innovation; $\theta(1)$ is the sum of the moving-average coefficients. Harris writes $V_t$ for value and $\varepsilon_t$ for its change. An alternative notation is $P_t=M_t+u_t$; Hasbrouck's $u_t$ is public news in the generalized Roll model, so this entry uses $s_t$.

\qffield{Conditions} Fundamental volatility comes only from \emph{unexpected} changes in the factors that set values. Predictable price changes that pay holders for carrying costs and risk, such as a zero-coupon bond's creep towards par, are not volatility; an unexpected change in rates is \citep[chap.~``Volatility'', pp.~411--412]{alpha-harris}. When news is common knowledge prices can change without trading; when few know it, prices move on high volume as the informed trade; and the price revisions that make up the adverse-selection spread component contribute to fundamental volatility \citep[same chap., p.~411]{alpha-harris}.

Transitory volatility includes bid--ask bounce, produced by the transaction-cost (transitory) spread component, and the price effects of large orders and cumulative uninformed imbalances that reverse when value traders and arbitrageurs act; it includes both the price change and its reversal. Value traders, arbitrageurs and dealers do not cause it but resolve it \citep[same chap., pp.~413--414]{alpha-harris}. Roll's model is the special case $s_t=\tfrac12SQ_t$ with trade signs independent of $w_t$, which gives the two-term variance $\sigma^2+\tfrac12S^2$ (see \hyperref[res:roll-spread-estimator]{Roll's estimator}).

\qffield{Practice} Fundamental factors differ by instrument: supply and demand, interest rates and storage costs for commodities; rates and credit quality for bonds; management, resources and technologies, product and input markets and rates for stocks; inflation, macroeconomic policy and trade and capital flows for currencies. Volatility is higher for commodities that are costly to store or perishable, especially with long production pipelines, variable demand or inelastic demand; for firms whose value depends on unresolved uncertainties, such as research outcomes or the growth behind high price--earnings ratios; for instruments exposed to political risk; and for the equity of highly leveraged firms \citep[same chap., pp.~411--413]{alpha-harris}. Because transitory price changes revert, negative serial correlation of price changes is a strong indicator of transitory volatility: over successive trades for bid--ask bounce, and over minutes to months for uninformed order imbalances \citep[same chap., p.~414]{alpha-harris}.

A variance ratio $V_{M,N}=[\operatorname{Var}(p_t-p_{t-M})/M]/[\operatorname{Var}(p_t-p_{t-N})/N]$ equals one for a random walk; when microstructure effects dominate short horizons, $V_{M,N}>1$ for $M<N$ \citep[sec.~8.7, eq.~(8.14), pp.~74--75]{s6-hasbrouck}. Transitory volatility moves with uninformed traders' transaction costs and is small in liquid markets \citep[chap.~``Volatility'', p.~414]{alpha-harris}. Regulators can have no lasting effect on fundamental volatility but can change transitory volatility, in either direction \citep[same chap., p.~410]{alpha-harris}. Circuit breakers slow fundamental price changes; their net effect on transitory volatility is indeterminate, lowering it when they restrain uninformed traders, order anticipators and bluffers and raising it when they restrain informed traders, dealers and arbitrageurs, and empirical studies have been inconclusive. Margins and transaction taxes cannot discriminate between these trader types \citep[chap.~``Bubbles, Crashes, and Circuit Breakers'', secs.~28.3.2--28.3.4, pp.~575--577]{alpha-harris}.

\qffield{Pitfalls} Negative serial correlation is not proof of transitory volatility: successive spot prices of a perishable commodity are prices of different items, for different delivery dates, so their negative correlation can be fundamental \citep[chap.~``Volatility'', box, p.~413]{alpha-harris}. Its absence is not proof of a random walk either: lagged price adjustment makes price changes positively autocorrelated \citep[sec.~8.6, Exercise~8.2, p.~74]{s6-hasbrouck}, and negative autocorrelation at some lags can offset positive autocorrelation at others to give a variance ratio near one \citep[sec.~8.7, p.~75]{s6-hasbrouck}. Roll's variance components understate transitory volatility when reversals take longer than one trade \citep[chap.~``Volatility'', p.~416]{alpha-harris}.

The split of $\operatorname{Var}(\Delta p_t)=\sigma_w^2+\operatorname{Var}(\Delta s_t)+2\operatorname{Cov}(w_t,\Delta s_t)$ beyond $\sigma_w^2$ is not identified; assigning the covariance term to transitory volatility is a convention, and Harris's ``sum of the two'' holds when the components are uncorrelated (Interpretation). Episodic and continuing volatility are not further sources. Harris uses episodic volatility for episodes of large price changes in short intervals \citep[same chap., p.~410]{alpha-harris}, a description rather than a third source; continuing volatility does not appear in the source.

\qffield{See also} \hyperref[res:roll-spread-estimator]{Roll's estimator} (special case and estimator); \hyperref[res:spread-component-identification]{component identification} (the same split for spreads: the generalized Roll model); \hyperref[meth:garman-klass-volatility]{Garman--Klass volatility} (measuring total volatility from ranges); \hyperref[mech:bubble-crash-dynamics]{bubbles and crashes} (episodes mixing both components); \hyperref[res:spread-determinants]{spread determinants} (volatility as a determinant); \hyperref[def:fundamental-value]{fundamental value} (Harris's value, conditioned on all information, not only public); \hyperref[def:intraday-overnight-returns]{returns by time scale}.

\end{qfdefinition}

\begin{qfmethod}[Garman--Klass volatility from open, high, low and close prices]\label{meth:garman-klass-volatility}
\qffield{Statement} Lehalle and Laruelle estimate intraday volatility over a fixed interval from its open, high, low and close prices $p^O,p^H,p^L,p^C$ and its average price $\bar p$ \citep[chap.~2, sec.~2.5.1, p.~172]{alpha-laruelle}:
\[ \mathrm{Vol}_{GK}=\frac{\sqrt{\tfrac12\left(p^H-p^L\right)^2-(2\log2-1)\left(p^C-p^O\right)^2}}{\bar p}. \]
\begin{enumerate}
\item Choose the interval. Lehalle and Laruelle use 5 minutes to show intraday seasonality and 15 minutes to compare markets \citep[same chap., secs.~2.2 and~2.5.1, pp.~146 and~172]{alpha-laruelle}.
\item Record the open, high, low and close of each interval and apply the formula.
\item Average across intervals or days; the comparison of FTSE~100 and IBEX~35 components used a 20-day moving average of daily averages \citep[same chap., sec.~2.5.1, p.~172]{alpha-laruelle}.
\end{enumerate}

\qffield{Notation} The superscripts $O,H,L,C$ avoid a clash with the transaction-cost component $C$. In log-ratio form,
\[
\sigma_{GK}=\sqrt{\tfrac12(\log H/L)^2-(2\log2-1)(\log C/O)^2}.
\]
Since $\log(p^H/p^L)\approx(p^H-p^L)/\bar p$ for small ranges, the two forms agree to first order (Interpretation).

\qffield{Conditions} Timeless (method); the applications use 2007--2016 European data. Volatility is not observed: each estimate presumes a model. Under an arithmetic Brownian price the natural estimator sums squared trade-to-trade changes, but on high-frequency data that estimator is polluted by bid--ask bounce, and with jumps they must be removed first. Garman--Klass is a measure not polluted too much by the bounce. Implied volatility measures something else, the uncertainty priced into options \citep[same chap., sec.~2.2, pp.~146--147]{alpha-laruelle}.

\qffield{Practice} Lehalle and Laruelle use it for the intraday seasonality of volatility, with volume, spread and book size, on Royal Dutch Shell, 2011--2016 (Figure~2.12, pp.~146--148), and to compare the intraday volatility of FTSE~100 and IBEX~35 components after MiFID, finding no significant effect of the UK's fragmentation on intraday volatility \citep[same chap., sec.~2.5.1, pp.~172--174]{alpha-laruelle}.

\qffield{Pitfalls} The coefficient $2\log2-1$ and the estimator's properties come from Garman and Klass (1980); Lehalle and Laruelle give the formula without derivation or assumptions. The estimate is per interval and relative to price; scaling it to a daily or annual figure needs an assumption about how variance adds across intervals (Interpretation). When the high and low are trade prices they include part of the spread, so on short intervals the range still carries some bounce (Interpretation).

\qffield{See also} \hyperref[def:volatility-components]{fundamental and transitory volatility} (what the total contains); \hyperref[meth:spread-volatility-regression]{spread--volatility regression} (a use of interval volatilities); \hyperref[res:roll-spread-estimator]{Roll's estimator} (the bounce that trade-based estimators pick up); \hyperref[def:intraday-overnight-returns]{returns by time scale} (open-to-close and overnight returns).

\qffield{Further reading} The original estimator \citep{garman-klass-1980}.

\end{qfmethod}

\begin{qfdefinition}[Intraday, overnight and close-to-close returns: volatility and correlation by time scale]\label{def:intraday-overnight-returns}
\qffield{Statement} For day $t$ with opening price $p^O_t$ and closing price $p^C_t$, Lehalle and Laruelle define the \emph{intraday} and \emph{overnight} returns \citep[chap.~2, sec.~2.5.1, eqs.~(2.2)--(2.3), p.~175]{alpha-laruelle}:
\[ r^{\mathrm{intra}}_t=\log p^C_t-\log p^O_t, \qquad r^{\mathrm{on}}_t=\log p^O_t-\log p^C_{t-1}, \]
and compute the correlation between two stocks separately on each series, over a six-month sliding window, averaging the pairwise correlations across an index's components \citep[same sec., pp.~174--175]{alpha-laruelle}. Interpretation: the close-to-close (\emph{extraday}) return is exactly their sum, $r^{\mathrm{cc}}_t=r^{\mathrm{on}}_t+r^{\mathrm{intra}}_t$. Volatility likewise has an intraday and an extraday scale: investors with extraday constraints bear extraday volatility, and intraday volatility only through its effect on that scale \citep[same chap., sec.~2.5, p.~169]{alpha-laruelle}.

\qffield{Notation} Superscripts $O$ and $C$ for open and close, as in \hyperref[meth:garman-klass-volatility]{Garman--Klass}, avoid a clash with the transaction-cost component $C$.

\qffield{Conditions} The two returns come from different market mechanisms: overnight the market is organised as an auction fixing, during the day as a continuous double auction. The authors hold that the opening fixing is less contaminated by microstructure noise than the closing fixing, and so better reflects investors' consensus on the price \citep[same chap., sec.~2.5.1, p.~177]{alpha-laruelle}. Interpretation: $\operatorname{Var}(r^{\mathrm{cc}})=\operatorname{Var}(r^{\mathrm{on}})+\operatorname{Var}(r^{\mathrm{intra}})+2\operatorname{Cov}(r^{\mathrm{on}},r^{\mathrm{intra}})$, so extraday risk mixes the two scales and their covariance.

\qffield{Practice} On the components of the CLES~60, DAX, CAC~40 and IBEX~35, 1991--2012, correlation rises mainly on the overnight side, and overnight correlation explains well the spikes in extraday correlation, to which it is closer; intraday correlation acts as a baseline that overnight correlation keeps bouncing above \citep[same sec., Figures~2.27--2.28, pp.~174--177]{alpha-laruelle}. Intraday correlation spiked in summer 2010 and in August 2011, and for three of the indices was in October 2010 at its highest level since 1991 \citep[same sec., pp.~174 and~176]{alpha-laruelle}. The authors read intraday correlation as an intermediate marker in a contagion process by which microscopic causes have macroscopic effects \citep[same sec., p.~174]{alpha-laruelle}.

\qffield{Pitfalls} The authors attribute part of the March--October 2010 rise in extraday correlation to rising intraday correlation, the consequence of a less noisy PFP, and offer growing high-frequency activity as one possible explanation \citep[same sec., p.~177]{alpha-laruelle}: a conjecture, not a test. Daily correlation spikes can have different mechanisms at different horizons; the authors find the spikes mainly overnight. An interpretation links overnight correlation to auction and news revaluation, and overnight volatility to information arrival. The auction is in the source; news arriving overnight is an interpretation. A principle said intraday frictions matter for daily risk when they change the closing or auction price or the path used by benchmark investors.

The source asks whether changes in intraday price formation show at the extraday scale \citep[same chap., sec.~2.5, p.~169]{alpha-laruelle}, and compares the noise in the two fixings (p.~177); benchmark paths are not discussed there. Interpretation: the returns depend on which prices are used, auction prints or last trades, and the open and close of a fragmented market may differ by venue.

\qffield{See also} \hyperref[meth:garman-klass-volatility]{Garman--Klass volatility} (volatility within intraday intervals; the FTSE~100 and IBEX~35 comparison of the same section); \hyperref[def:volatility-components]{fundamental and transitory volatility} (variance across horizons); \hyperref[def:price-formation-process]{price formation process}.

\end{qfdefinition}

\begin{qfmethod}[Spread--volatility regression in levels and logs]\label{meth:spread-volatility-regression}
\qffield{Statement} Lehalle and Laruelle test how the spread moves with volatility in five steps \citep[sec.~A.10, pp.~292--298]{alpha-laruelle}:
\begin{enumerate}
\item Compute the spread $\psi_i$ and the volatility $\sigma_i$ over the same intervals; their example uses 15-minute intervals of Rio Tinto from January 2012.
\item Fit $\psi_i=a_0+a_1\sigma_i+\epsilon_i$ by ordinary least squares: $\hat a_1=\sum_i(\sigma_i-\bar\sigma)(\psi_i-\bar\psi)/\sum_i(\sigma_i-\bar\sigma)^2$, $\hat a_0=\bar\psi-\hat a_1\bar\sigma$.
\item Inspect the distributions: both variables looked log-normal.
\item Fit $\log\psi_i=a_0+a_1\log\sigma_i+\epsilon_i$, which in levels is the multiplicative model $\psi=A_0\sigma^{a_1}e^{\epsilon}$, $A_0=e^{a_0}$.
\item If $\epsilon\sim N(0,\sigma_\epsilon^2)$, then $\mathbb{E}[\psi]=A_0\sigma^{a_1}e^{\sigma_\epsilon^2/2}$, so predict levels with $\hat\psi=\hat A_0\sigma^{\hat a_1}e^{\hat\sigma_\epsilon^2/2}$, which they call the unbiased predictor.
\end{enumerate}
A positive $\hat a_1$ is consistent with the spread rewarding market makers for the risk they take: the higher the volatility, the more they ask \citep[chap.~2, sec.~2.3.1, p.~150]{alpha-laruelle}.

\qffield{Notation} $\psi$ is the spread, $\sigma$ the volatility over the interval, $a_0,a_1$ the intercept and slope, $\sigma_\epsilon^2$ the residual variance. $A_0$ is the level intercept, not the adverse-selection component $A$.

\qffield{Conditions} As of the 2008--2016 data used; method timeless. The least-squares assumptions are those of the simple linear model: the regressor measured without error; errors with mean zero, constant variance and no correlation with each other or with the regressor; normal errors for the $t$ tests \citep[app.~A, sec.~A.10.1, pp.~293--295]{alpha-laruelle}. Across small-tick instruments, Lehalle and Laruelle report the relation $\psi_i=a+b\,\sigma_i/\sqrt{N_i}$, with $N_i$ the average number of trades per day and $a,b$ shared within an asset class; for large-tick instruments, $\psi_i=a+(b_1\sigma_i+b_2\,\eta\,\mathrm{tick})/\sqrt{N_i}$, where $\eta=N_c/(2N_a)$ is half the ratio of price continuations to alternations, smaller the more the price bounces \citep[chap.~2, sec.~2.2, pp.~144--145]{alpha-laruelle}.

\qffield{Practice} For Rio Tinto, the level fit gave $\hat a_0=2.4894$, $\hat a_1=0.0485$, $R^2=0.2404$; the log fit $\hat a_0=0.5609$, $\hat a_1=0.2252$, $R^2=0.1553$; the back-transformed predictor, with $\hat\sigma_\epsilon=0.2509$, explains 0.1097 of the variance of the spread in levels \citep[app.~A, Figures~A.14, A.17 and~A.18, pp.~296--298]{alpha-laruelle}. Daily spreads rise with daily volatility on BNP Paribas and BP, 2008--2012; points off the line in August 2011 coincide with the short-selling ban, when market makers could not meet their commitments in financial stocks \citep[chap.~2, sec.~2.3.1, Figure~2.13, pp.~150--151]{alpha-laruelle}.

\emph{Fit quality}. The total variability splits into explained and residual parts, and $R^2=\mathrm{EV}/\mathrm{TV}=1-\mathrm{RV}/\mathrm{TV}$, equal in the simple model to the squared correlation $r^2$ (sec.~A.10.1, p.~295). The equality of the two forms holds for least squares with an intercept; for the back-transformed predictor of step~5 they differ, and Lehalle and Laruelle's $R'^2$ (0.1097) uses the explained form $\sum(\hat\psi_i-\bar\psi)^2/\sum(\psi_i-\bar\psi)^2$ (p.~298): on simulated data the two forms give 0.2984 and 0.2984 for least squares and 0.2990 against 0.2709 for such a predictor.

\emph{The averaging effect} \citep[app.~A, sec.~A.14.2, pp.~319--321]{alpha-laruelle}. For an index $I_t=\sum_kw_kS^k_t$ the spread aggregates, $\psi^I=\sum_kw_k\psi^k$, ``but this relation is not true for the volatility''; so the coefficients of the spread--volatility regression on index quantities need not equal the weighted averages of the stock regressions' coefficients.

On the DAX, April to July 2012, the index regression gave $\hat a^I_0=0.0283$ and $\hat a^I_1=2.6641$ against weighted averages 0.0277 and 2.4734, higher ``due to the index volatility which is not the weighted average of the stock volatilities'' (Figure~A.28, p.~321).

\qffield{Pitfalls} The fit is descriptive: in the example volatility explains a quarter or less of the spread's variance. Unless the comparison controls for them, other variables drive both series. Harris's general warning applies: observable determinants are correlated with each other and with more than one primary factor, so tests need controls \citep[chap.~``Bid/Ask Spreads'', p.~314]{alpha-harris}.

Tick size: when the spread is stuck at one tick, liquidity providers add size at the best quotes instead, so in large-tick assets volatility may show in depth rather than in the spread \citep[chap.~2, sec.~2.2, p.~148]{alpha-laruelle}. Time of day: spread, volatility, volume and book size share an intraday seasonality, high after the open, so pooling intervals across the day mixes the seasonal pattern with the relation \citep[same sec., pp.~145--148]{alpha-laruelle}. Trading activity: the cross-sectional relation scales volatility by $1/\sqrt{N_i}$ \citep[same sec., pp.~144--145]{alpha-laruelle}.

The source does not discuss venue and depth as regression controls. Causation runs both ways, since a wider spread adds bounce to trade-based volatility (Interpretation). The predictor is unbiased only with known $a_0$, $a_1$ and $\sigma_\epsilon^2$; with estimates plugged in it is approximately unbiased (Interpretation). The underlying mechanism is that market makers widen spreads when volatility raises adverse-selection and inventory risk, and the relation is cleaner when volatility is scaled by trading activity and the tick is accounted for; Lehalle and Laruelle summarise that ``High volume combined with low volatility and a tight bid--ask spread characterize a liquid stock'' \citep[chap.~2, sec.~2.3, p.~149]{alpha-laruelle}.

The spread identity holds for spreads in price units, since the index's best ask and bid are the weighted sums of the components'; for spreads in basis points the index spread is a price-weighted, not a $w$-weighted, average (3.96 against 4.73~bp on simulated components). The text defines $\bar a_1=\sum_kw_k\hat a^k_0$ (p.~320), a printing slip for $\sum_kw_k\hat a^k_1$.

Index volatility in price units is at most $\sum_kw_k\sigma^k$ for nonnegative weights, with equality only under perfect correlation (1.598 against 2.376 at correlation 0.4). The index regression's slope can therefore be steeper than the stock slopes' average even when every stock follows the same spread--volatility relation: in a simulation with a common relation of slope 2.5 the stock slopes average 2.507 and the index slope is 3.108. If the index volatility were a fixed fraction $c<1$ of $\sum_kw_k\sigma^k$, the index slope would be the common slope divided by $c$ (derived).

\qffield{See also} \hyperref[res:spread-determinants]{spread determinants} (the theory being tested: volatility as a primary determinant); \hyperref[meth:garman-klass-volatility]{Garman--Klass volatility} (an interval volatility); \hyperref[def:spread-measures]{spread measures} (which spread to regress).

\end{qfmethod}

\begin{qfmechanism}[Bubbles and crashes: how mispricing builds and unwinds]\label{mech:bubble-crash-dynamics}
\qffield{Statement} A \emph{bubble} is a rise of prices substantially above fundamental values, quickly or over a long period; a \emph{crash} is a very fast fall, often but not always after a bubble. Since fundamental values are not common knowledge, bubbles and crashes can occur at all \citep[chap.~``Bubbles, Crashes, and Circuit Breakers'', sec.~28.1, p.~556]{alpha-harris}. The typical sequence \citep[same chap., sec.~28.1.1, pp.~556--557]{alpha-harris}:
\begin{enumerate}
\item Buyers grow overly optimistic about new technologies or markets, often without recognising what prices already reflect, and push prices up.
\item Momentum traders and order anticipators buy and accelerate the rise; early buyers grow confident and early sellers buy back to stop losses.
\item Value traders and arbitrageurs cannot sell enough (no position to sell, too little capital to carry a short, short selling hard) or will not (they expect uninformed buying to continue, or doubt their own value estimates).
\item Sellers start trading aggressively: long-term holders, early buyers taking gains, contrarians, value traders, arbitrageurs. Momentum buyers lose interest and optimists lose confidence.
\item Late buyers sell to stop losses, margin calls and stop orders force more selling, and order anticipators sell ahead of them: prices fall quickly. Panic can take prices below value, and a bounce follows.
\end{enumerate}
\emph{The price accelerator} (Harris, after a manuscript by Treynor): rising prices move wealth from pessimists with short positions to optimists with long ones; the richer optimists may buy more, and the losing pessimists may buy to cover margin calls, so prices rise further; the sellers are mild optimists and mild pessimists \citep[same chap., box, p.~556]{alpha-harris}.

\qffield{Conditions} Timeless (mechanism); the case material is historical, to 2002. Most bubbles and crashes involve both fundamental and transitory volatility. A bubble usually starts on good news, a fundamental change; pseudo-informed traders then overreact, the transitory part. It usually bursts on bad news, but the volume and price response is out of proportion to the news because of the overpricing that preceded it; the reversion towards value counts as transitory volatility, since it occurs only because uninformed traders inflated prices, and panic below value adds more \citep[same chap., sec.~28.1.2, pp.~557--558]{alpha-harris}.

\qffield{Practice} Harris's 1987 account shows the mechanisms together: portfolio insurance traded like stop orders, and uncertainty about how much of it was waiting kept buyers out; NYSE order printers could not keep up, so traders did not know whether their orders had filled; index futures led the cash market down and traded at a discount of more than 10 percent once index arbitrageurs stopped; and buy-side traders could not reach Nasdaq dealers \citep[same chap., sec.~28.2.2, pp.~559--563]{alpha-harris}. Crashes rarely have a single cause: traders become confused about risk, attend to the wrong issues and, when events move fast, panic \citep[same chap., sec.~28.2.8, p.~571]{alpha-harris}.

\qffield{Pitfalls} A large fall without news of commensurate importance does not show that the market failed. Most crashes correct earlier pricing errors; if they were failures, prices would rebound fully once traders saw them below value, but after most crashes prices rebound little \citep[same chap., box, p.~559]{alpha-harris}. Attribution therefore looks at the reversal after the event, at what preceded it, and at overloaded systems, forced selling and broken links between markets, as in 1987 (Interpretation, from the cases on pp.~557--563). Bubbles are hard to identify in real time, and analysts generally understand the causes of extreme volatility better after the fact \citep[same chap., p.~555]{alpha-harris}. Among the stories that start a bubble, Harris's general account names technology and new markets, and his Japanese example turns on loose monetary policy \citep[same chap., sec.~28.2.7, pp.~570--571]{alpha-harris}.

\qffield{See also} \hyperref[def:volatility-components]{fundamental and transitory volatility} (the two components in each phase); \hyperref[def:outside-spread]{outside spread} (why value traders wait); \hyperref[res:spread-determinants]{spread determinants}. See also \hyperref[mod:depth-overflow-flash-crash]{the toy model of the Flash Crash} and \hyperref[ev:may-2010-flash-crash]{the May 2010 Flash Crash}.


\end{qfmechanism}
